
The empirical forcefield (FF) provides us with the most trivial conceptual understanding of the PES, since it seeks  an analytical expression for the total energy, where, in the most general form, it can be written as a series expansion of the  coordinates ($\bm{r}$) for two, three, up to $n$ interacting atoms in the system:
\begin{align}\label{eq:total_energy_atom_pos_expand}
E(\bm{R})=V_{0}+&\sum_{i}V^{(1)}(\bm{r}_{i},\bm{\beta}^{(1)})
               +\frac{1}{2!}\sum_{i,j}V^{(2)}(\bm{r}_{i},\bm{r}_{j},\bm{\beta}^{(2)})+\cdots\nonumber\\
               &+\frac{1}{n!}\sum_{i,j,\ldots,n}V^{(n)}(\bm{r}_{i},\bm{r}_{j},\ldots,\bm{r}_{n},\bm{\beta}^{(n)}).
\end{align} 
Here, each term of $V^{(i)}$ takes a well-defined mathematical form that is formulated based on the known physics of (classical) mechanics and electrostatics. With this, both the total system energy and its derivates with respect to the atomic positions can be easily and efficiently calculated, enabling us to study the structures and properties of materials at the length  (more than a million atoms) and time (up to milliseconds\cite{pierce2012routine}) scales that are relevant to experimental observations. Examples that are specifically relevant to material science include: the determination of phonon density-of-states (see \cref{sect:phonon}) to supplement the inelastic neutron scattering experiments,\cite{cheng2019simulation} thermal transport,\cite{esfarjani2011heat} crystallisation and polymorphism in molecular crystals,\cite{raiteri2005exploring} phase transformation,\cite{pols2021atomistic} and many others. We would like to highlight that, while the current discussions  around \cref{eq:total_energy_atom_pos_expand} focus on atom-atom interaction, the idea of many-body expansion is general across many aspects of material modelling. For example, in determining the total interaction energy of non-covalently interacting molecular aggregates using the fragmented approach,\cite{mahadevi2016cooperativity,wen2012practical} the series will be expanded based on the positions of $n$--interacting molecules, in which the higher--order terms are particularly important for the interaction energies contributed by the long-ranged effect of charge polarisation. As will be further discussed in \cref{sec:anharmonic_vibs}, the anharmonic phonon vibrations in solids are primarily dictated by the non-negligible higher-order terms $(n>2)$ in   \cref{eq:total_energy_atom_pos_expand}. 

The key idea to review the constructions of major  forcefields used in material and molecular simulation is to be compared with the later review in the development of machine-learning forcefield, which learns the PES $E(\bm{R})$ solely from the training data without explicitly assuming the underlying functional forms that are necessary for describing these interactions.

\subsubsection{Parameterisation of the Forcefield}

FF methods are commonly referred to as being \emph{empirical} because they contain a set of parameters $\{\bm{\beta}^{(i)}\}$ [\cref{eq:total_energy_atom_pos_expand}] that is specifically fitted to a particular material system for a particular functional form of the FF, such that the materials' properties obtained from the FF  $\mathcal{P}^{FF}(\{\bm{\beta}^{(i)}\})$ reproduce the targeted properties $\mathcal{P}^{target}$ obtained from either experiments or higher level theoretical calculations. In another word, the values of $\{\bm{\beta}^{(i)}\}$  are determined by minimising the following error:
\begin{equation}\label{eq:FF_fitting}
\min_{\bm{\beta}}\left\Vert \mathcal{P}^{FF}(\{\bm{\beta}^{(i)}\})-\mathcal{P}^{target}\right\Vert.
\end{equation}
Before we move on to discuss some of the common FF types in materials modelling, we will provide a brief account of different choices $\mathcal{P}^{target}$ that have been used in fitting the FF parameters.

For the use of experimental data as $\mathcal{P}^{target}$, several common choices have been reported in the literature. For instances, the measured sublimation enthalpy of crystal can be directly correlated with the cohesive energy calculated from a model FF. These thermochemical data, which are often augmented with structural data such as lattice constants, have been applied to fit the FF parameters for crystalline materials, such as molecular crystals.\cite{williams1967nonbonded,warshel1970consistent,pyzer2016optimized,bunte2000molecular} %This can be further supplemented with experimental lattice constants to ensure the FF can also correctly reproduce the experimental structures. 
For biological and aqueous systems, neutron scattering is capable of providing detailed structural and dynamic informations, particularly for its high sensitivity to the atomic position of the hydrogen atoms. In this regard, the energy and momentum dependent neutron scattering cross section $S(\bm{Q},E)$  has been applied as $\mathcal{P}^{target}$ in refining the FF  to obtain more accurate structures for biomolecules\cite{lynch2017automated,pluhackova2015extension} and the hydration environments of the physiologically important ions.\cite{martinek2018calcium,kohagen2016accounting,duboue2017hydration} The vibrational energy levels, which is determined by the Hessian of the PES function [\cref{eq:force_constant}], can be compared to the vibrational frequencies measured from the Raman spectroscopy to refine $\{\bm{\beta}^{(i)}\}$. Such strategy had been applied in developing the FF for $n$-butane,\cite{murphy1991harmonic} and $\beta$-\ce{Ga2O3}.\cite{dohy1982raman} The Raman intensities, which can be  calculated from the dipole-dipole autocorrelation functions extracted from molecular dynamic trajectories, have also been applied as $\mathcal{P}^{target}$ to refine the polarisable FF for water.\cite{sidler2018efficient} Finally, X-ray and electron diffraction measurements can provide direct and detailed information on the equilibrium bond lengths, bond angles and van der Waals (vdW) interaction distances, which are some of the important parameters needed in the FF.\cite{sun1998compass,zgarbova2013toward,aytenfisu2017revised} 

By parameterising the FF directly against experimental data provides a significant advantage that enables the predicted properties $\mathcal{P}^{FF}$ for new systems modelled with the FF to become directly comparable with the experiments for verifying the theoretical predictions. However, the qualities and the consistency of the experimental data, the account of thermal effects, as well as the difficulties in simultaneously fitting multiple inter-dependent parameters, make the parameterisation of the empirical FF a non-trivial, often  tedious process. Some of these challenges may be mitigated with the usage of training data that is generated from \emph{ab initio} calculations.

For examples, by simultaneously fitting both the intra- and inter-molecular PES to the one that is consistently derived from periodic DFT calculations, Neumann\cite{neumann2008tailor} developed a recipe for constructing tailored-made FF for crystals of flexible organic molecules, thus enabling the accurate ranking of the relative stabilities of different polymorphs of organic crystals which are typically separated by energy gaps in the order of $\sim 1$ kJ/mol. Achieving such a level of accuracy has important implications in   guiding the screening of crystallisations of drug molecules in the pharmaceutical industry.\cite{neumann2015combined} Note that Neumann adapted a simple Buckingham\cite{buckingham1938classical} potential of the form $E(r)=A\exp(Br)-C/r^6$ in his FF, which is a variation of the Lennard-Jones potential commonly used in many empirical forcefields (see Fig. 2 in \citet{heinz2013thermodynamically}),  to describe the vdW interaction.  One of the challenges here is to reach a balanced fitting for the sharp-rising exponential function at short-distance and slowly varying $1/r^6$-dependency at long range. The application of Symmetry-Adapted Perturbation Theory (SAPT)\cite{misquitta2005intermolecular}, which is an electronic structure theoretical method, enables the intermolecular interaction energies to be decomposed into  contributing terms, whereby each term can be individually fitted to a well-defined analytical function,\cite{misquitta2016ab} thus mitigating the problem of cross-coupling among the FF parameters. In an effort to make the FF parameterisation become more of a `black-box' method, Grimme proposed a quantum mechanically derived forcefield (QMDFF) for organic molecules that can be parameterised using only the results from quantum mechanical calculations,\cite{grimme2014general} particularly the atomic charges and the Hessian, albeit the resulting FF is only aimed to reproduce the PES close to a specific equilibrium structure obtained from quantum mechanical calculation. 
%\textcolor{blue}{More examples of parameterising the FF against the \emph{ab initio} data needed, especially beyond molecular system.}

Having briefly reviewed the strategies for parameterising the empirical FFs, we moved on to give a high-level overview of the functional forms of the major FFs that are commonly used in atomistic modelings. The discussion here will mainly focus on how different types of  FFs are categorised according to the fundamental physics behind the atom-atom interactions that are aimed to be captured by the specific functional forms of the FFs, with some modifications following the work from \citet{heinz2013thermodynamically} More specifically, they are (a) common harmonic that describes interactions among the bonded (including both covalent/ionic and vdW) atoms, (b)  many-body/reactive FFs in which the intramolecularly bonded atoms are allowed to vary during the simulation, and finally, (c) polarisable FFs which are able to account for the change in the strong electrostatic interactions resulted from the deformations of the electron densities as the interacting atoms approach each others. 

\subsubsection{Common Harmonic Terms}
In real practice, most of the FFs applied in material modelling do not use a functional form  that is dependent directly on the Cartesian coordinates of the interacting atoms [\emph{c.f.} \cref{eq:total_energy_atom_pos_expand}], but the internal coordinates that represent the bond lengths ($r_{ij}$, between a pair of atoms), bond angles ($\theta_{ijk}$, among a triplet of atoms) and torsional angles ($\varphi_{ijkl}$, among a quartet of atoms). Such a choice is not only based on our intuitive understanding on the material chemistry, but also, as has been demonstrated in many prior works,\cite{baker1999geometry,buvcko2005geometry,echenique2006definition} that the molecular motions along the curvilinear space spanned by (an appropriate choice for the set of) internal coordinates tend to follow low-energy-barrier pathways along the PES. Consequently, this can improve the stabilities in gradient-based geometry optimisations, as well as the efficiencies in sampling the PES  for structure predictions (see \cref{sect:MonteCarlo}). 

With this, the lowest level of approximation to construct the FF is to assume all bonded interactions are \emph{harmonic}, whereas the electrostatic interactions between pairs of atoms are modelled in the form of simple Coulomb interactions between \textit{fixed charges}. For example, the CHARMM FF that is commonly used in modelling biomolecular systems takes the following specific form:\cite{mackerell1998all}
\begin{align}\label{eq:charmm_general}
    E=&\sum_{ij \mbox{\scriptsize{,bonded}}}V_{\mbox{\scriptsize{bond}}}\left(r_{ij}\right)
    +\sum_{ijk,\mbox{\scriptsize{bonded}}}V_{\mbox{\scriptsize{angle}}}\left(\theta_{ijk}\right)+\sum_{ijkl,\mbox{\scriptsize{bonded}}}V_{\mbox{\scriptsize{proper-tor}}}\left(\varphi_{ijkl}\right)\nonumber\\
&+\sum_{ijkl,\mbox{\scriptsize{bonded}}}V_{\mbox{\scriptsize{improper-tor}}}\left(\chi_{ijkl}\right)+\sum_{ij,\mbox{\scriptsize{nonbonded}}}V_{\mbox{\scriptsize{Coulomb}}}\left(r_{ij}\right)+\sum_{ij,\mbox{\scriptsize{nonbonded}}}V_{\mbox{\scriptsize{LJ}}}\left(r_{ij}\right),
\end{align}
where each term in the above expression is
\begin{subequations}
\begin{align}
V_{\mbox{\scriptsize{bond}}}\left(r_{ij}\right)&=K_{r,ij}\left(r_{ij}-r_{ij,\mbox{\scriptsize{eq}}}\right)^2\label{eq:bond}\\
V_{\mbox{\scriptsize{angle}}}\left(\theta_{ijk}\right)&=K_{\theta,ijk}\left(\theta_{ijk}-\theta_{ijk,\mbox{\scriptsize{eq}}}\right)^2\label{eq:angle}\\
V_{\mbox{\scriptsize{proper-tor}}}\left(\varphi_{ijkl}\right)&=\frac{1}{2}K_{\varphi,ijkl}\left[1+\cos\left(n\varphi_{ijkl}-\varphi_{ijkl,\mbox{\scriptsize{eq}}}\right)\right]\label{eq:proper_tor}\\
V_{\mbox{\scriptsize{improper-tor}}}\left(\chi_{ijkl}\right)&=K_{\chi,ijkl}\left(\chi_{ijkl}-\chi_{ijkl,\mbox{\scriptsize{eq}}}\right)^2\label{eq:improper-tor}\\
V_{\mbox{\scriptsize{Coulomb}}}\left(r_{ij}\right)&=\frac{1}{4\pi\varepsilon}\frac{Q_{i} Q_{j}}{r_{ij}}\label{eq:Coulomb}\\
V_{\mbox{\scriptsize{LJ}}}\left(r_{ij}\right)&=\varepsilon_{ij}\left[\left(\frac{r_{ij,\mbox{\scriptsize{eq}}}}{r_{ij}}\right)^{12}-2\left(\frac{r_{ij,\mbox{\scriptsize{eq}}}}{r_{ij}}\right)^{6}\right]\label{eq:LJ}.
\end{align}
\end{subequations}
FFs adapting this type of functional from are usually dubbed as the Class I FF, which, apart from the CHARMM FF, also include FFs such as AMBER,\cite{pearlman1995amber} OPLS-AA,\cite{jorgensen1996development} DREIDING\cite{mayo1990dreiding} and many others. The so-called Class II FF (\emph{e.g.} the COMPASS\cite{sun1998compass} FF) includes cubic and quartic terms that accounts for the anharmonic effects (see \cref{sec:anharmonic_vibs}) in molecular motions, as well as the cross-coupling terms, among the bonded atoms. The latter terms are included because  internal coordinates  are not necessarily orthogonal to each other, changing the torsional angles in a molecule, for example, may simultaneously requires certain bond angles to be changed at the same time in order to avoid the steric clashes of the bulky side groups. The inclusion of the cross-coupling term increases the flexibility of the FF to capture this important physical behaviour, which can be  important for dynamic simulations to guide the molecule to transform from one conformation to another by traversing through part of the PES with lower transition barriers. 

The simple functional forms adopted in both the Class I and II FF to describe the PES also imposes some limitations in modelling the physical and chemical behaviours of materials. Here, we focuses on the two major limitations.

Firstly, the FF assumes and maintains a fixed bond topology, meaning no chemical reactions (forming/breaking of the covalent bonds) can occur throughout the simulation. Fundamentally, this is limited by the functional form [\emph{c.f.} \cref{eq:bond}] chosen to describe the bond stretching, which does not include the long-range asymptotic behaviour $E(r\to\infty)=0$ that promotes bond breakage. In addition, \cref{eq:bond} only contains \emph{one} well-defined local minimum at the (parameterised) equilibrium bond distance between the two atoms $i$ and $j$ that are predefined to be bonded together, meaning this particular functional form of the FF is physically constraining the motion of the atoms in the vicinity around the local minimum at  $r_{ij,\mbox{\scriptsize{eq}}}$ on the PES. Secondly, both Class I and II FFs model the electrostatic interactions [\cref{eq:Coulomb}] between two atoms by using fixed point charges on atoms. Whilst this simple form of Coulomb interaction is extremely fast to evaluate in large-scale atomistic modelling, it ignore the important physics of the electronic charge densities around atoms, the shape of which can change when two atoms come close to each other, contributing to the important polarisation effects that consequently affect the strength of the Coulomb interactions between two interacting atomic moieties. For instances, ignoring the charge polarisation may lead to 40 \% error in the interaction energies between two water molecules when the intermolecular oxygen--oxygen distances fall under 3 $\mbox{\AA}$.\cite{ferenczy2001modeling} In order to model the PES  over a wider range and with better accuracy,  reactive and polarisable FFs have been proposed and developed to overcome these two limitations, respectively, which will be briefly reviewed in the following sections.

\subsubsection{Many-Body and Reactive Forcefields}

A useful approach to overcome the limitation imposed by the predefined bond topology/connectivity is to consider the total system energy as a sum of  the energies from individual atoms in the system, each interacting with multiple neighbouring atom that can be identified as those fall within a predefined cut-off distance $r_c$ to the central atom. This  constitutes a many-body FF as the number of atoms that can be considered as interacting with the central atom grows as $r_c$ increases, which is particularly useful for developing FFs for solids,  as the atoms in solids typically have  higher coordination numbers compared to the number of bonds that can form with the C atoms in organic molecules. It also allows the number of neighbouring atoms to changes, thus making it possible to model the change in the total system energy as a result of a chemical reaction. As shall be demonstrated in the latter sections [\cref{sec:physics_informed_representation}], the fundamental concept behind many-body FFs has played a particularly significant role in the design of atom-centred structure descriptors for modern machine-learning methods in material modelling. For this reason, and its importance in modelling inorganic crystalline solids, which is of the main interest for the present review, we will focus on three major many-body/reactive FFs in more details, namely the embedded atom method (EAM), the bond-order potential (BOP) and the charge optimised many-body (COMB) potential.
 
\paragraph{\textbf{The embedded atom method}} was first proposed by \citet{daw1984embedded} for modelling the impurities, defects and surfaces of metals, to overcome the unsatisfactory results in using pairwise FF in treating systems such as H in the metal lattice. The EAM is inspired from the Hohenberg-Kohn formulation of the DFT\cite{hohenberg1964inhomogeneous} which views each atom in the solid as an impurity that is embedded in the potential generated by all the other atoms in the solid. With this, the total energy in the EAM takes the following form:
\begin{equation}\label{eq:EAM_total_energy}
E=\sum_{i}^N F(\rho_{h,i})+\frac{1}{2}\sum_{i\neq j}\phi(r_{ij}).
\end{equation}
Here, $F(\rho_{h,i})$ is the energy required to embed atom at the site $i$ of the host lattice, which is a function of the electron density of all atoms in the host lattice at this site $\rho_{h,i}$. The latter can be built by summing a (model) electron densities at all sites except $i$, \emph{i.e.} $\rho_{h,i}=\sum_{j\neq i}\rho_j^a(R_{ij})$. The  core-core repulsive interaction is described by the term $\phi(r_{ij})$. Similar to the pair potentials discussed in the previous section, practical application of EAM potential will require a specific functional form for the embedding and core-core repulsion potentials,  where some of the popular choices are briefly discussed below which a specific focus on the physical motivations behind the chosen functional forms.

In the original formulation of the EAM for Ni and Pd,\cite{daw1984embedded}  the  charge densities $\rho_i(r)$ were obtained from the parameters for the atomic wavefunctions for the $s$ and $d$ orbitals determined from the \emph{ab initio} Hartree-Fock calculations.\cite{clementi1974roothaan} With this, the total charge density of an atom $\rho(r)=N_s\rho_s(r)+(N-N_s)\rho_d (r)$ where $N$ is the total number of outer electrons and $N_s$ is the number of $s$ electrons.  From here, $F(\rho_{h,i})$ was obtained as a cubic spline function satisfying the constraint of it being of a monotonically decreasing function from $F(\rho)\vert_{\rho=0}=0$, which is set to be the zero point energy. Because  $N_s$ cannot be determined from homonuclear data, it becomes a free-parameters in the potential which must be determined by fitting to the experimental data (such as lattice constants, elastic constants and cohesive energies of the equilibrium structure), in order to determine the exact shape of $F(\rho)$. For the core-core repulsion energy, a simple form of Coulombic interaction $\phi_{ij}(r)=Z_i(r)Z_j(r)/r$ is adopted. The distance-dependent effective charge $Z(r)$ is also numerically fitted, which must satisfy the boundary conditions where $Z(0)=Z^{*}$ (the atomic number) and must asymptotically decrease to 0 when $r\to a_0$, which is the atomic radius. In a subsequent work, \citet{foiles1986embedded} developed a consistent form of the embedding function and short-ranged repulsive pair interactions for Cu, Ag, Au, Ni, Pd, Pt and their binary alloys. In this work, the information about $F(\rho)$ was obtained from the knowledge of system energies away from $\rho_{eq}$, the charge densities at the equilibrium crystal structure, when the latter is expanded or shrinked. In particular, this takes the form of $E(a)=-E_{sub}(1+\Delta a)e^{-\Delta a}$, where $E_{sub}$ is the sublimation enthalpy of the metal, and $\Delta a=\left[\frac{a}{a_{0}}-1\right]/\sqrt{\frac{E_{sub}}{9B\Omega}}$ measures the deviation from the equilibrium lattice constant $a_{0}$ where $B$ is the Bulk modulus and $\Omega$is the equilibrium cell volume per atom. 

Later in the work by \citet{ackland2004development}, where EAM potential was developed for modelling P impurity embedded in Fe, the embedding potential was chosen to take the simple form of $F(x)=-\sqrt{x}+a_2 x^2+a_4 x^4$. The $\sqrt{x}$-dependency takes the inspiration from the tight-binding model\cite{finnis1984simple} with the quadratic and quartic terms accounts for the kinetic energy of electrons in the hopping and free-electron regime. The model parameters were obtained by fitting to the data derived from \emph{ab initio} calculations, including crystal structures, surface energies, lattice constants, as well as vacancy and defect formation energies. In the other works which aimed at modelling the Ni-H\cite{angelo1995trapping,baskes1997trapping} and W-H\cite{wang2017embedded} systems, the function $F(x)=c_6 x^2+c_7 x+\frac{c_8 x^{5/3}}{c_4+x}$ was proposed to be a suitable form for modelling H in the metal matrix. The spherical atomic densities was assumed which was further augmented with a exponentially decaying cut-off function $f_{cut}(r)=\exp\left[\frac{1}{r-r_{cut}}\right]$. The usage of cut-off function with similar form is quite common in the construction of atom-centred structural descriptors\cite{bartok2013representing,behler2011neural} for atomistic machine-learnings. Finally, \citet{sheng2011highly} developed a set of highly optimised EAM potentials for 14 FCC metals. To ensure the properties modelled by the parameterised EAM potentials, the \emph{ab initio} determined lattice constants and energies were further re-scaled to the experimental values during the parameterisation process. They adopted a quintic spline function for the embedding function, which enables them to obtain a smooth interpolation of the Hessian matrix, leading to the successful reproduction of the phonon dispersion relationship for Au that is closer to the 80 K experimental results from neutron diffraction, compared with that derived from pure \emph{ab initio} calculation.  By explicitly including the \emph{ab initio} molecular dynamics data for simulating the melting and liquid states of the metals, the EAM potential is capable of sampling the liquid structures and dynamics for the metals that have been parameterised. 

\paragraph{The bond-order potential} which can be represented by the Tersoff potential\cite{tersoff1986new} that is designed for covalent materials, such as silicon, to model its surface reconstruction behaviours that cannot be satisfactorily described by a simple two-body and three-body potential due to a fundamental change in the atomic coordination environments upon the formation of covalent surfaces. Similar to the two body potentials, the Tersoff potential model the total system energy as the sum of energies for all interacting atomic pairs $E=\sum_{j>i}V_{ij}$, where the pair interaction energy is the balance between the attractive and repulsive interaction energies that both take an exponential form of distance-dependency\cite{abell1985empirical}:
\begin{equation}\label{eq:tersoff}
V(r_{ij})=f_c(r_{ij})\left[A^{\mbox{\scriptsize{repul}}}\exp\left(-\lambda^{\mbox{\scriptsize{repul}}}r_{ij}\right)
-B^{\mbox{\scriptsize{attr}}}_{ij}\exp\left(-\lambda^{\mbox{\scriptsize{attr}}}r_{ij}\right)\right],
\end{equation}
in which $A$, $B_{ij}$ and $\lambda$'s are parameters and $f_c(r_{ij})$ is the smooth cut-off function, thus the summation $\sum_{j>i}$ operation will be performed for all atom pairs that fall within the cut-off radius $r_{cut}$ of any given central atom. Unlike the simple pair-potentials, the attractive parameter $B_{ij}$ takes into account the bond orders, as well as the change in the bonding strength due to the change in the local bonding environments (\emph{i.e.} the bond lengths and angles formed between atom $i$ and other atoms $k\neq j$). More explicitly, 
\begin{subequations}
    \begin{align}
    B_{ij}&=B_{0}\exp(z_{ij}/b)\label{eq:tersoff_b}\\
        z_{ij}&=\sum_{k\neq i,j}\left(\frac{w(r_{ik})}{w(r_{ij})}\right)^n \frac{1}{c+\exp(-d\cos\theta_{ijk})}\label{eq:tersoff_z}\\
        w(r_{ij})&=f_c(r_{ij})\exp\left(\lambda^{\mbox{\scriptsize{attr}}}r_{ij}\right)\label{eq:tersoff_w}
    \end{align}
\end{subequations}
 Physically, $z_{ij}$ measures the number of competing bonds with bond $ij$ and the parameter $b$ determines how rapidly the strength of bond $ij$ falls with the increase in the bond coordination number. $w(r_{ij})$  is the unscaled bond strength, and the first term in \cref{eq:tersoff_z} is the ratio in bond strengths between bond $ik$ and $ij$, where the power $n$ determines how significant bonds between closely neighbouring atoms are weighted against the distanced atoms. Finally, the second term in \cref{eq:tersoff_z} gives the weighting factor based on the bond angle $\theta_{ijk}$. Fundamentally, a change in the local coordination number changes the shape of the charge densities around the central atom, thus affects the overlap of charge densities between the neighbouring atoms which consequentially changes the bond strengths.\cite{abell1985empirical} In a later development, for example, \citet{brenner1990empirical} extended the bond order potential to account for the occurrence of chemical reactions by further refining the $B$ parameters depending on the nature of the chemical bonds (\emph{e.g.} the $\sigma$ and $\pi$ bonds). Similar idea of local coordination dependency is also reflected in Grimme's --D3\cite{grimme2010consistent}  scheme to correct the inaccuracy of DFT in dealing with the non-covalent interactions, where the elemental $C_6$ coefficient for the dispersion interaction can be attenuated depending on the specific coordination number of that element in a particular compound. Many structural fingerprints [\emph{c.f.} \cref{sec:feature_eng}] that are devised for atomistic machine learnings, such as SOAP\cite{bartok2013representing} and Behler's atom-centred symmetry functions,\cite{behler2011atom} can see their inspirations from the construction of bond-order potentials, such as the one of Tersoff.
 
Several material systems have been parameterised with their corresponding Tersoff potential, mostly against \emph{ab initio} training data. For instances, \citet{lindsay2010optimized} optimised the potential for graphene and carbon nanotubes against their experimental phononic properties, with a specific aim in applying the potential for modelling the phonon transport properties. More specifically, in order to model the bonding environments in the carbon nanotube, an additional term that depends on the  torsion angles coupling four atoms in the system is included for the $B$ parameter. \citet{munetoh2007interatomic} parameterised the potential for the Si-O system, which is capable of reproducing the properties of various crystalline silica polymorphs as well as its amorphous state. The parameters optimised by \citet{kumar2016optimised} for  B-N enable them to perform molecular dynamic studies for examining the fracture behaviours of two-dimensional BN nanosheets in atomistic details. \citet{billeter2006ab}  developed the Tersoff potential for the silicon oxynitride system, which can be applied for investigating its deposition behaviour on the silicon substrate, for its promising physical properties to be applied as gate dielectrics in the microelectronics. Similar potential has been developed for a series of semiconductor nitrides,\cite{karaaslan2020assessment} and In-doped GaAs systems.\cite{titantah2007bond}
 
\paragraph{\textbf{Charge optimised many-body potential} } In all types of many-body potentials discussed above, the electrostatic energy between two interacting atoms is assumed to possess a fixed charge that is independent from its chemical environment. This poses a fundamental limitation in using the potential to model systems that involves complex chemical bonding environments that may also be dynamically varying,\cite{liang2013classical} such as the heterogeneous interfaces formed between two different materials, as well as  the chemical reactions occurring at these interfaces where the charges on atoms changes as they migrate from one material from another. 

A remedy for this problem is to extend the QEq charge equilibration approach\cite{rappe1991charge} that is widely used in the \textit{ab initio} atomistic simulations into empirical potentials, which leads to the development of the charge optimised many-body (COMB) potentials. In this case, the functional form for the total energy of the system is not only just the function of the atomic position vector $\bm{r}$, but also the vector $\bm{q}$ of the point charges on individual atom. With this,  the COMB potential takes the following form:\cite{yu2007charge}
\begin{equation*}
    E(\bm{q},\bm{r})=V_{\mbox{\scriptsize{short}}}(\bm{q},\bm{r}) +V_{\mbox{\scriptsize{es}}}(\bm{q},\bm{r})+V_{\mbox{\scriptsize{vdw}}}(\bm{r})+V_{\mbox{\scriptsize{corr}}}(\bm{r}).
\end{equation*}
Here, the short-ranged interaction $V_{\mbox{\scriptsize{short}}}$ takes the form of the bond-order potential [\cref{eq:tersoff}], except now both the attractive and repulsive terms have extra term that is charge dependent.  $V_{\mbox{\scriptsize{vdw}}}$ describes the van der Waals interaction between two atoms and takes the form of a 12-6 potential. $V_{\mbox{\scriptsize{corr}}}$ is introduced in the latest version of the COMB potential to additionally improve the description of bond angle terms. Finally, the electrostatic interactions \begin{equation*}
V_{\mbox{\scriptsize{es}}}=V_{\mbox{\scriptsize{self}}}+V_{\mbox{\scriptsize{qq}}}+V_{\mbox{\scriptsize{qZ}}}+V_{\mbox{\scriptsize{polar}}}
\end{equation*} 
is the summed interaction energies contributed from: (a) the self energy $V_{\mbox{\scriptsize{self}}}(\bm{q},\bm{r})$ describes the energy cost associated with the formation of a charge ion, which generally takes the form of a Taylor series expansion based on the atomic charge, (b) $V_{\mbox{\scriptsize{qq}}}(\bm{q},\bm{r})=\frac{1}{2}\bm{q}\mathbf{J}^{qq}\bm{q}$ is the Coulomb interaction between two spherical $s$--type charge densities with the strength of the interaction given by the tensor $\mathbf{J}^{qq}$, similarly
(c) $V_{\mbox{\scriptsize{qZ}}}(\bm{q},\bm{r})=\frac{1}{2}\bm{q}\mathbf{J}^{qZ}\bm{Z}$ is the Coulomb interactions between electrons and the nuclei. Finally, to further improve the accuracy of the electrostatic interactions, the polarisability interactions (d) $V_{\mbox{\scriptsize{polar}}}=\sum_i \frac{\mu_i^2}{2P_i}+\sum_{i}\bm{\mu}_i\cdot \bm{E}^{q}_i+\sum_{i,j}\bm{\mu}_i\mathbf{T}_{ij}\bm{\mu}_j$ is included where $\bm{\mu}_i$ is the electronic dipole moment,  $\bm{E}^{q}_i$ is the electric field and $\mathbf{T}_{ij}$ is the dipole-dipole interaction tensor (c.f. \citet{stone2013theory}). 

Since the atomic charges are no longer  fixed parameters but variables in the COMB potential, in the dynamic simulation of material systems using this potential, the effect of charge flows amongst the atomic species are additionally included in the Lagrangian to optimised the charge distribution in the simulated system, that is further subjected to the condition of electronegativity equalisation.\cite{rappe1991charge,rick1994dynamical}  In this regard, the COMB potential can also be considered as a type of polarisable forcefield (discussed in detailed in the next section), except the key difference here is that the COMB potential is an reactive potential where the bonded atoms are allowed to change during the simulation.

Within this framework of modelling the atom-atom interactions, several technologically important materials, the interfacial behaviours in which are critical for microelectronic devices, have been successfully parameterised and modelled using the COMB potential. Some of the notable examples include Si/\ce{SiO2},\cite{yu2007charge} amorphous silica,\cite{shan2010second} Hf/\ce{HfO2},\cite{shan2010charge} Ti/\ce{TiO2},\cite{cheng2014charge} as well as \ce{Al2O3}.\cite{choudhary2015charge}


\subsubsection{Polarisable Forcefields}\label{sect:polarisable_FF}

\begin{figure*}[thb]
    \centering
    \includegraphics[width=\textwidth]{cooperativity_illustration.png}
    \caption{\pol{(a-b) Measurements of electron density shift (EDS) in pairs of non-covalently bounded molecules.\cite{iribarren2021evaluation} Here, the EDS is defined as $\mbox{EDS}=\rho(abc\ldots n)-\sum_{i=1}^{N}\rho(i)$, which is the difference in the electron density for a mutually interacting cluster of $N$ molecules, and the sum of electron densities for isolated monomers in their respective geometries exactly as in the cluster. Green and magenta represent positive and negative values for the EDS, respectively. (a) Evolution of the percentage of EDS inside the cutoff (a.u.) and EDS values (in-plots) with respect to the cutoff for \ce{H2O}$\cdots$\ce{H2O}. (b) EDS maps at the 0.001 a.u. cutoff for trimers of malonaldehyde arranged in linear and zigzag shapes. (c)  Benchmark results for 3-body energies $\Delta E_{3b}$ for 69 small molecular trimers bounded by non-covalent interactions, which were extracted from 23 different molecular crystals.\cite{rezac2015benchmark} Here the total $\Delta E_{3b}$ is broken down into the contributions from the electronic polarisations and dispersions to showcase the relationships between them. (d) Convergence of the lattice energy for the benzene crystal calculated at the highly accurate coupled cluster level of theory using many-body expansion approach,\cite{yang2014ab} which shows that including the contributions from up to the 4-body term is necessary to converge the lattice energy to the experimental value (adopted from \citet{beran2014new}). (e) Erratic convergence of the many-body energies (see \cref{eq:p-body-nonadd})  towards the full cluster energies for  four different (\ce{H2O})$_{16}$ clusters.\cite{ouyang2014trouble}}}
    \label{fig:cooperativity_illustration}
\end{figure*}


\pol{\paragraph{\textbf{Preliminary:  non-additive interactions in many-body systems}} For the majority of the forcefield models discussed above, the total interaction energy of the entire system is generally assumed to be a simple sum of the interaction energies amongst all subsets of bodies in the system [\emph{c.f.} \cref{eq:total_energy_atom_pos_expand}]. In other words, given an $N$ body system, once parameterised, the interaction energy $V^{(n)}(\{\bm{r}_{n}\},\bm{\beta}^{(n)})$ for a subset of $n\subseteq N$ bodies is independent of the presence of the rest of bodies $\{i\in N:i\notin n\}$ in the system. However, this is not necessarily true even for simple material systems. Fundamentally, as illustrated in \cref{fig:cooperativity_illustration}(a-b), this is originated from the polarisability of the electron density centred on each body, where upon mutual contact, the electron densities can polarise each other through space in a geometry-dependent way, causing changes in the overall electron density distributions in the entire system, even they are bounded with (weak) non-covalent interactions. Since the interactions among electron densities play significant contributions to the total interaction energy, it can be appreciated that the total system energy of a many-body system will no longer be a simple summation of the  energies for all interacting pairs, each consider in isolation from the other bodies. This is because the latter energy is modified when the pair is brought into the contact with the remaining bodies in the whole system.}


\pol{To formalise this argument,\cite{kaplan1995non} consider the total energy $E(abc\cdots n)$ for a system with $N$ bodies. In a non-interacting system, the total energy is a sum of individual one-body energies, \emph{i.e.} $E(abc\cdots n)=E_{1b}(abc\cdots n)=\sum_{i=1}^{N}E(i)$. In an additive system, the total system energy needs to take into account only the contributions from all the two-body interactions in the system
$E(abc\cdots n)=E_{1b}(abc\cdots n)+\Delta E_{2b}(abc\cdots n)$, in which
\begin{equation}\label{eq:two-body-nonadd}
    \Delta E_{2b}(abc\cdots n)=\sum_{i<j}^{{N \choose 2}}E(ij)-a_{2N}^{1}E_{1b}(abc\cdots n),
\end{equation}
where $a_{mn}^{k}=\frac{(n-k)!}{(n-m)!(m-k)!}$ and $E(ij)$ is the total energy of a two-body subsystem in isolation. As discussed, in a non-additive system, the two-body terms $\Delta E_{2b}(abc\cdots n)$ must be further corrected by the presence of the other bodies that are also interacting with the two-bodies, such that the total system energy can be written as the following series:
\begin{equation}\label{eq:nonadd-total-energy}
E(abc\cdots n)=E_{1b}(abc\cdots n)+\sum_{p=2}^{N}\Delta E_{pb}(abc\cdots n),
\end{equation}
which can be considered as a variant of \cref{eq:total_energy_atom_pos_expand}, and it can be shown by induction that the non-additive $p$-body $(p\leq N)$ interaction energy
\begin{equation}\label{eq:p-body-nonadd}
    \Delta E_{pb}(abc\cdots n)=\sum_{i<j<\ldots<p}^{{N \choose p}}E(ij\cdots p)-\sum_{i=1}^{p-1} a_{pN}^{i}\Delta E_{pb}(abc\cdots n).
\end{equation}
}

\pol{It is clear from \cref{eq:p-body-nonadd} that,  calculating the non-additive $p$-body energies for an $N$-body system in general faces similar  sampling challenges in high-entropy systems (\cref{sect:HEA}), as the number of different $p$-body configurations grows factorially as a function of $p$ and $N$. Tackling this challenge, specifically in the field of non-covalent interactions in (supra)molecular systems.\cite{mahadevi2016cooperativity} Generally, clusters of polar molecules exhibit large non-additive many-body energies. For instance, benchmark calculations [\cref{fig:cooperativity_illustration}(c)] shows that the average $\Delta E_{3b}$ is $\sim 0.46$ kcal/mol for water trimer,\cite{rezac2015benchmark} whereby the largest $\Delta E_{3b}$  for water trimer was found to be $-2.42$ kcal/mol, which contributes significantly towards its stability (18.4$\%$ of the total interaction energy) that is comparable in magnitudes to the usual two-body terms. For non-polar  molecules that are predominantly bound together via dispersion interactions, such as benzene, in which the mutual interactions between benzene molecules will not significantly redistribute the charge density in the monomers, the corresponding $\Delta E_{3b}$'s in the benzene trimers are usually small. This has been exploited in an earlier study\cite{yang2014ab} to utilise the many-body expansion [\cref{eq:nonadd-total-energy}] for accurately determining the lattice energy of benzene crystal using quantum-chemistry method based on localised atomic orbitals, where sufficiently good agreement with experimental result is already reachable by truncating \cref{eq:nonadd-total-energy} at the 4-body term [\cref{fig:cooperativity_illustration}(d)]. This is not generalisable to polar molecules, however, as shown in \cref{fig:cooperativity_illustration}(e), where $\Delta E_{nb}$ shows strongly oscillating behaviours for three model examples of water clusters,\cite{ouyang2014trouble} making the expansion [\cref{eq:nonadd-total-energy}] difficult to converge. Various approaches have been developed to overcome such a challenge, and it is beyond the scope of discussion here, thus we direct interested readers to other reviews in this field.\cite{beran2016modeling}  }

\pol{\paragraph{\textbf{Popular polarisable forcefields in material modellings}} The need to capture the atomic position dependent deformation in the electronic density  (\emph{i.e.} electron polarisations) for accurate modelling the non-additive interactions in atomistic simulation transcends into the more complex functional forms and parameterisations that needs to be added into the atom-atom potential, which leads to the (specialised) development of polarisable forcefields. The COMB potential discussed in the previous section is a variant of he polarisable forcefield, but this term is more commonly used for categorising the forcefields that have been better developed for (bio)molecular simulations. A possible reason behind is that the magnitude of the polarisation components can be accurately quantified by decomposing it from the total interaction energy between two non-covalently bounded molecular moieties using SAPT,\cite{misquitta2005intermolecular,misquitta2016ab} enabling it to be efficiently parameterised with a well-defined  functional form with solid physical origins. Therefore, the primary focus of this section is to provide a brief overview on the three major strategies for including polarisation effects in forcefields for molecular simulations.\cite{jing2019polarizable} }

\pol{In the forcefields such as AMOEBA\cite{ponder2010current} and AMBER,\cite{cieplak2001molecular,wang2012development} electronic polarisation is accounted with the \emph{induced dipole model}. Here, distributed multipole analysis\cite{stone2013theory,cisneros2012application} is first applied to expand the molecular charge densities into atom-centred multipoles up to  quadrupolar level for obtaining an accurate description of molecular electrostatics. On each polarisable site, an  induced dipole in response to the external field $\bm{\mu}_{\mbox{\scriptsize{ind}}}=\alpha \bm{E}_{\mbox{\scriptsize{ext}}}$ is introduced, in which $\alpha$ is the electronic polarisability. Within this framework, the additional contribution to the total electrostatic interaction energy in the system is modelled as $E_{\mbox{\scriptsize{ind}}}=\sum_{i}\frac{1}{2}\alpha_i^{-1}\bm{\mu}^2_{i,\mbox{\scriptsize{ind}}}$, where the summation is performed over all polarisable sites. The molecular polarisability is determined iteratively using Thole's damped interaction method,\cite{thole1981molecular} which considers the  all polarisable sites can mutually interact with each other, which changes the charge distribution in the system until it becomes stabilised.}

\pol{On the other hand, the \emph{Drude oscillators model}\cite{lemkul2016empirical} takes a more classical  approach to account for the polarisation effects using point charges, which, has been implemented in forcefield such as CHARMM-Drude.\cite{baker2011development,jana2015charmm} In the Drude model, electronic polarisation is modelled as a point charge attached to the polarisble site with a parameterised force constant $k_i$, such that the induced electrostatic interaction takes the form of a usual harmonic potential $E_{\mbox{\scriptsize{ind}}}=\frac{1}{2}\sum_i  k_i\bm{d}_i^2$, where $d_i$ is the displacement of the point charge.}

\pol{Lastly, the \emph{fluctuating charge model}, which is used in, for example, the AMBER-$\sigma\pi$ forcefield,\cite{chen2007qtpie,patel2004charmm} takes similar approach at the COMB potential discussed above to optimise the atomic charges $\{ q_i \}$ in the system, where the total system energy is minimised by allowing charges to flow from one atom to another while maintaining the overall charge balance in the entire system. From here, the induced electrostatic energy $E_{\mbox{\scriptsize{ind}}}=\sum_i\left(\chi_i q_i+\eta_i q_i^2\right)$. Using the atomic ionisation potential (IP) and electron affinity (EA), the Mulliken electronegativity is then $\chi_i=\frac{1}{2}(\mbox{IP}+\mbox{EA})=-\mu_i$, whereas the Parr-Pearson hardness can be determined as $\eta_i=\frac{1}{2}(\mbox{IP}-\mbox{EA})$.}

\pol{\paragraph{\textbf{Exemplary applications of polarisable forcefields in condensed matter modellings}} Although less common, some investigations have been dedicated to parameterise and apply the polarisable forcefields in modelling technologically important materials in the condensed phases. \citet{rathnayake2020evaluation} parameterised the AMOEBA forcefield specifically for metal halide perovskites against the \emph{ab initio} reference data, although molecular dynamic simulations performed using the as-developed forcefields lead to an overestimation of the phase transition temperatures in these materials. It highlights the challenges in accurately modelling the electrostatic interactions between localised (molecular cations) and delocalised (inorganic frameworks) electronic states in a hybrid system, that is an over-lasting challenge in material modellings. Similarly, polarisation forcefield has been applied to study metal-organic frameworks (MOFs) and their gas uptaking capacities,\cite{mcdaniel2015evaluation,becker2017polarizable,becker2018polarizable} which highlighting the importance of including the polarisation effects in accurate determination of the thermodynamics of gas absorption in MOF. Although neglecting such effect may also lead to acceptable result, the outcome is most likely an error cancellation effect arises from the overestimation of the dispersion energies. Another directly relevant field of application is surface chemistry for sensors and catalysts. In this regard, for examples, \citet{schyman2013exploring} parameterised the OPLS\cite{jorgensen1995importance}-AAP forcefield for water adsorption on carbon surfaces, and achieved good agreement with the thermochemical data obtained from high-level quantum chemical calculations. Polarisbale CHARMM forcefields for gold\cite{wright2013golp} and silver\cite{hughes2013biomolecular} have also been developed, enabling the efficient investigations of the interactions between biomolecules and metal nanoparticles for sensing applications using atomistic modellings.}

\pol{\citet{walz2018phase} developed polarisable forcefield to examine the phase transition and molten behaviours of alkaline halide crystals. Predicting the solubilities of molecular crystal is critically important for the development of new pharmaceutical drugs, \citet{schnieders2012structure} tested the applicability of the AMOEBA forcefield for modelling the structure, thermodynamics and the solubility of the hexanamide organic crystal, achieving an accuracy of 1.1 kcal/mol in predicting the solvation free energy. By parameterising against DFT results, \citet{corradini2015dft} developed forcefields that accurately predicted the stabilities for the polymorphs of \ce{TiO2} and \ce{SiO2}, and subsequently applied it to model their amorphous structure, achieving excellent agreement with the experimental neutron structural factors.  }
